\documentclass[11pt,letterpaper]{article}
\usepackage[margin=1in]{geometry}
\usepackage{amsmath,amssymb,amsthm}
\usepackage{physics}
\usepackage{hyperref}
\usepackage{cleveref}
\usepackage{booktabs}
\usepackage{longtable}
\usepackage{array}
\usepackage{multirow}

\newtheorem{theorem}{Theorem}[section]
\newtheorem{proposition}[theorem]{Proposition}

\numberwithin{equation}{section}
\renewcommand{\theequation}{S11.\arabic{equation}}

\title{Supplement S11: The Complete $(2,3,5,7)$ Connection Table \\
\large Supporting Manuscript \S15}
\author{UDT Collaboration}
\date{}

\begin{document}
\maketitle

\begin{abstract}
We present the complete table showing how every derived quantity in UDT
traces to the four multiplicities $(2,3,5,7)$ arising from the unique
Diophantine triple $(j=1/2,\, l=1,\, |\kappa_\mathrm{max}|=3)$.
From one metric and one Diophantine identity, four numbers generate:
particle masses, coupling constants, mixing angles, cosmological
parameters, and the CMB boundary condition. This supplement serves as the
master index for the derivation chain: one metric $\to$ one Diophantine
triple $\to$ four numbers $\to$ all physics.
\end{abstract}

\tableofcontents

%----------------------------------------------------------------------
\section{The four numbers}
\label{sec:four_numbers}
%----------------------------------------------------------------------

The Diophantine identity (Supplement~S4)
\begin{equation}
  (2j+1)^2(2l+1)(2|\kappa_\mathrm{max}|+1) = \binom{2l+2|\kappa_\mathrm{max}|+1}{2l+1}
\end{equation}
selects $j = 1/2$, $l = 1$, $|\kappa_\mathrm{max}| = 3$ uniquely, producing
four multiplicities:

\begin{equation}
\label{eq:four_numbers}
\boxed{
\begin{array}{cccc}
  2j+1 & 2l+1 & 2|\kappa_\mathrm{max}|-1 & 2|\kappa_\mathrm{max}|+1 \\
  \mathbf{2} & \mathbf{3} & \mathbf{5} & \mathbf{7}
\end{array}
}
\end{equation}

These are the first four primes. The base space dimension is
$n = 2l + 2|\kappa_\mathrm{max}| + 1 = 9$.

%----------------------------------------------------------------------
\section{Derived combinations}
\label{sec:combinations}
%----------------------------------------------------------------------

All physically relevant quantities are built from products, powers, and
binomial coefficients of $(2, 3, 5, 7)$:

\begin{center}
\begin{tabular}{lcl}
\toprule
Expression & Value & Role \\
\midrule
$(2j+1)^2$ & 4 & Spin-squared factor \\
$(2l+1)^2$ & 9 & Orbital-squared factor \\
$(2j+1)(2l+1)$ & 6 & Proton coefficient \\
$(2j+1)^2(2l+1)$ & 12 & --- \\
$(2j+1)^2(2|\kappa_\mathrm{max}|-1)/(2l+1)$ & 20/3 & Muon coefficient \\
$(2j+1)^2(2l+1)^2$ & 36 & $\binom{9}{2}$, $\alpha$ numerator \\
$(2j+1)^2(2l+1)(2|\kappa_\mathrm{max}|+1)$ & 84 & $\binom{9}{3}$, pion coefficient \\
$(2l+1)^2(2|\kappa_\mathrm{max}|-1)$ & 45 & $\theta_{13}$ denominator \\
$(2j+1)^2 + (2l+1)^2$ & 13 & $\theta_{12}$ denominator \\
$(2|\kappa_\mathrm{max}|-1)/(2j+1)^2$ & 5/4 & p-n splitting factor \\
\bottomrule
\end{tabular}
\end{center}

%----------------------------------------------------------------------
\section{Master connection table: Masses}
\label{sec:masses}
%----------------------------------------------------------------------

\begin{longtable}{p{2.5cm}p{5cm}p{3.5cm}c}
\toprule
Particle & Formula & Numbers used & Supplement \\
\midrule
\endfirsthead
\toprule
Particle & Formula & Numbers used & Supplement \\
\midrule
\endhead
\textbf{Electron} & $m_e$ (anchor) & --- & S3 \\
\textbf{Pion} & $m_\pi = \binom{9}{3}\pi\,m_e = 84\pi\,m_e$ & $2 \cdot 3 \cdot 7 \cdot \pi$ & S3 \\
\textbf{Muon} & $m_\mu = \frac{20\pi^3}{3}\,m_e$ & $\frac{2^2 \cdot 5}{3}\cdot\pi^3$ & S3 \\
\textbf{Proton} & $m_p = 6\pi^5\,m_e$ & $2 \cdot 3 \cdot \pi^5$ & S3 \\
\textbf{Neutron} & $m_n = m_p + \frac{5}{4}\alpha C$ & $\frac{5}{2^2} \cdot \alpha \cdot C$ & S5 \\
\textbf{Neutrino} & $m_\nu = \frac{\alpha^3 m_e}{4}$ & $\frac{1}{2^2}\cdot\alpha^3$ & S6 \\
\textbf{Mesons} & $m = E_n(\kappa) \cdot C$ & $\kappa \in [-3,+3]$, $C$ & S2 \\
\textbf{Baryons} & $m = E_n(\kappa) \cdot C$ on $[0, r_*/\varphi]$ & $\kappa$, $C$, $\varphi$ & S2 \\
\bottomrule
\end{longtable}

where $C = 4\pi^2 m_e r_*$ is the calibration constant and
$\varphi = (1+\sqrt{5})/2$ is the golden ratio (baryon sub-cavity).

%----------------------------------------------------------------------
\section{Master connection table: Couplings}
\label{sec:couplings}
%----------------------------------------------------------------------

\begin{longtable}{p{3cm}p{5cm}p{3cm}c}
\toprule
Coupling & Formula & Numbers used & Supplement \\
\midrule
\endfirsthead
\toprule
Coupling & Formula & Numbers used & Supplement \\
\midrule
\endhead
\textbf{$1/\alpha_\mathrm{EM}$} & $36\pi/I_2$ & $2^2 \cdot 3^2 \cdot \pi$ & S5 \\
\textbf{Bridge $|\kappa|=1$} & $\mathcal{S}^2 E_1^2 = 1/18$ & $1/(2\cdot 3^2)$ & S8 \\
\textbf{Bridge $|\kappa|=2$} & $\mathcal{S}^2 E_1^2 = \pi^2/36$ & $\pi^2/(2^2\cdot 3^2)$ & S8 \\
\textbf{Bridge $|\kappa|=3$} & $\mathcal{S}^2 E_1^2 = 1/8$ & $1/(2^3)$ & S8 \\
\bottomrule
\end{longtable}

%----------------------------------------------------------------------
\section{Master connection table: Mixing angles}
\label{sec:mixing}
%----------------------------------------------------------------------

\begin{longtable}{p{2.5cm}p{5cm}p{3.5cm}c}
\toprule
Angle & Formula & Numbers used & Supplement \\
\midrule
\endfirsthead
\toprule
Angle & Formula & Numbers used & Supplement \\
\midrule
\endhead
$\sin^2\theta_{12}$ & $\frac{(2j+1)^2}{(2j+1)^2+(2l+1)^2} = \frac{4}{13}$ & $\frac{2^2}{2^2+3^2}$ & S7 \\
$\sin^2\theta_{23}$ & $\frac{(2j+1)^2}{2|\kappa_\mathrm{max}|+1} = \frac{4}{7}$ & $\frac{2^2}{7}$ & S7 \\
$\sin^2\theta_{13}$ & $\frac{1}{(2l+1)^2(2|\kappa_\mathrm{max}|-1)} = \frac{1}{45}$ & $\frac{1}{3^2 \cdot 5}$ & S7 \\
\bottomrule
\end{longtable}

%----------------------------------------------------------------------
\section{Master connection table: Cosmology}
\label{sec:cosmology}
%----------------------------------------------------------------------

\begin{longtable}{p{3cm}p{5cm}p{3cm}c}
\toprule
Quantity & Formula & Numbers used & Supplement \\
\midrule
\endfirsthead
\toprule
Quantity & Formula & Numbers used & Supplement \\
\midrule
\endhead
$\mu_g$ & $\pi\mu/13$ & $\pi/(2^2+3^2)$ & S10 \\
$\kappa_\mathrm{cosmo}$ & $(3/2)\mu_g$ & $3/2$ & S10 \\
$\beta_\mathrm{cosmo}$ & $-\cos(\pi/5)\,\mu_g^2$ & $\cos(\pi/5)$ & S10 \\
$\gamma_\mathrm{cosmo}$ & $(2/3)\mu_g^3$ & $2/3$ & S10 \\
$z_\mathrm{CMB}$ & $e^{\phi(r_*)}-1$ & (boundary value) & S10 \\
$c^2 = 2GM/r_*$ & Misner--Sharp closure & (self-consistency) & S10 \\
\bottomrule
\end{longtable}

%----------------------------------------------------------------------
\section{Master connection table: Exterior algebra}
\label{sec:exterior}
%----------------------------------------------------------------------

\begin{longtable}{p{3cm}p{5cm}p{3cm}c}
\toprule
Quantity & Formula & Numbers used & Supplement \\
\midrule
\endfirsthead
\toprule
Quantity & Formula & Numbers used & Supplement \\
\midrule
\endhead
Base space dim & $2l+2|\kappa_\mathrm{max}|+1 = 9$ & $3+7-1$ & S4 \\
$\dim\Lambda^2$ & $\binom{9}{2} = 36$ & $\alpha$ sector & S5 \\
$\dim\Lambda^3$ & $\binom{9}{3} = 84$ & Pion sector & S3 \\
$\dim\Lambda^4$ & $\binom{9}{4} = 126$ & Transition & S6 \\
$\dim\Lambda^6$ & $\binom{9}{6} = 84$ & Neutrino (Hodge) & S6 \\
Damping ratio & $84/126 = 2/3$ & Grade 3/4 & S6 \\
Even/odd ratio & $36/84 = 3/7$ & Grade 2/3 & S6 \\
Orbit matchings & $\binom{4}{3} = 4$ & $=(2j+1)^2$ & S9 \\
\bottomrule
\end{longtable}

%----------------------------------------------------------------------
\section{The derivation chain}
\label{sec:chain}
%----------------------------------------------------------------------

The complete derivation chain from the metric to all physics:

\begin{enumerate}
  \item \textbf{Metric}: $ds^2 = -e^{-2\phi}c^2dt^2 + e^{2\phi}dr^2 + r^2d\Omega^2$.

  \item \textbf{Dirac equation on the metric}: Tetrad $\to$ spin connection
    $\to$ Form-T system $\to$ angular separation.

  \item \textbf{Angular sector}: Spinor harmonics on $S^2$ with quantum
    numbers $(j, l, |\kappa_\mathrm{max}|)$.

  \item \textbf{Diophantine selection}: The self-consistency condition
    selects $(j, l, |\kappa_\mathrm{max}|) = (1/2, 1, 3)$ uniquely.

  \item \textbf{Four multiplicities}: $(2j+1, 2l+1, 2|\kappa_\mathrm{max}|-1, 2|\kappa_\mathrm{max}|+1) = (2, 3, 5, 7)$.

  \item \textbf{All physics}: From these four numbers and $\pi$, construct
    masses, couplings, mixing angles, and cosmology.
\end{enumerate}

\begin{proposition}[Completeness]
Every dimensionless quantity predicted by UDT can be expressed as a rational
function of $(2, 3, 5, 7, \pi)$ and the vacuum integral $I_2$.
\end{proposition}

The only non-algebraic input is $I_2 = \int e^{2\phi}\,dr$, which requires
solving the nonlinear ODE numerically. All other quantities are exact
algebraic expressions.

%----------------------------------------------------------------------
\section{Counting free parameters}
\label{sec:free_params}
%----------------------------------------------------------------------

\subsection{Input parameters}

UDT has the following inputs:
\begin{enumerate}
  \item $m_e = 0.51100$~MeV: electron mass (anchor/unit),
  \item $\phi_0 = -\cos(\pi/5)$: metric depth (derived from the Diophantine identity),
  \item $\mu^2 = \pi/3$: screening mass (derived from the metric),
  \item $r_* = 7 - 1/80$: cavity size (derived from the metric).
\end{enumerate}

Of these, only $m_e$ is a measured input (it sets the unit of mass).
The other three are geometric constants derived from the metric and
the Diophantine identity.

\subsection{Comparison with the Standard Model}

The Standard Model has $\sim 19$ free parameters (6 quark masses, 3 lepton
masses, 3 CKM angles + 1 phase, 3 gauge couplings, Higgs mass and VEV,
$\theta_\mathrm{QCD}$). UDT has 1 free parameter ($m_e$), with all others
derived from geometry.

\begin{center}
\begin{tabular}{lcc}
\toprule
 & Standard Model & UDT \\
\midrule
Free parameters & $\sim 19$ & 1 \\
Masses predicted & 0 & 6+ \\
Couplings predicted & 0 & $\alpha_\mathrm{EM}$ \\
Mixing angles predicted & 0 & 3 (PMNS) \\
\bottomrule
\end{tabular}
\end{center}

%----------------------------------------------------------------------
\section{What is NOT determined by $(2,3,5,7)$}
\label{sec:not_determined}
%----------------------------------------------------------------------

For completeness, the following quantities are \emph{not} purely determined
by the four multiplicities:
\begin{itemize}
  \item $I_2$: requires numerical ODE integration,
  \item Individual Dirac eigenvalues $E_n(\kappa)$: require numerical shooting,
  \item The cosmological $\phi(r)$ profile shape: requires numerical integration,
  \item Baryon sub-cavity eigenvalues: depend on $\varphi = (1+\sqrt{5})/2$
    (which is algebraic but not directly from $(2,3,5,7)$).
\end{itemize}

However, the \emph{ratios} and \emph{algebraic forms} of the eigenvalues
(e.g., $E_1(\kappa=-1) = 2\sqrt{2}/3$) are algebraic expressions involving
only small integers and square roots.

%----------------------------------------------------------------------
\section{Verification}
\label{sec:verification}
%----------------------------------------------------------------------

This table is the master index. Each entry references a specific supplement
(S1--S10, S12) and analysis script (01--10). The full verification chain:
\begin{itemize}
  \item \texttt{analysis/01\_vacuum\_profile.py} $\to$ $\phi(r)$, $I_2$
  \item \texttt{analysis/02\_eigenvalues.py} $\to$ $E_n(\kappa)$
  \item \texttt{analysis/03\_sources.py} $\to$ $\mathcal{S}(\kappa)$
  \item \texttt{analysis/04\_bridge\_identities.py} $\to$ $\mathcal{S}^2 E_1^2$
  \item \texttt{analysis/05\_alpha\_em.py} $\to$ $1/\alpha$, $m_n - m_p$
  \item \texttt{analysis/06\_particle\_masses.py} $\to$ meson/baryon masses
  \item \texttt{analysis/07\_angular\_sector.py} $\to$ $m_\pi$, $m_\mu$, $m_p$, Diophantine
  \item \texttt{analysis/10\_wedge\_product.py} $\to$ exterior algebra, Hodge duality
\end{itemize}

\end{document}
